\documentclass[11pt,onside]{article}
\usepackage{sectsty}
\usepackage{geometry}
\usepackage[utf8]{inputenc}
\usepackage[spanish,es-nodecimaldot,es-tabla ]{babel}
\usepackage[ruled,vlined,linesnumbered,lined,boxed,commentsnumbered,onelanguage]{algorithm2e}
\usepackage{bm}
\usepackage{upgreek}
\usepackage{amsmath}
\usepackage{amssymb}
\usepackage{microtype} 
\usepackage[svgnames]{xcolor} 
\usepackage{titlesec}
\usepackage{enumitem}
\usepackage{caption}
\usepackage{subcaption}
\usepackage{fancyhdr}
\usepackage{graphicx}
\usepackage{mwe}
\usepackage{cite}
\usepackage{makecell}
\usepackage{multirow}
\usepackage{xcolor,colortbl}
\usepackage{setspace}
\usepackage{caption}
\usepackage{lipsum}
\usepackage{float}
\usepackage{listings}
\geometry{
 letterpaper,
 left=20mm,
 right=20mm,
 tmargin=30mm,
 bmargin=30mm,
 }
\definecolor{Custom_color_1}{rgb}{0.56,0.83,0.29}
\definecolor{Custom_color_2}{rgb}{0.2,0.2,0.8}
\definecolor{mygreen}{RGB}{28,172,0} % color values Red, Green, Blue
\definecolor{mylilas}{RGB}{170,55,241}
\captionsetup{justification=raggedright,labelfont={color=Custom_color_2,bf}}
\setlength{\headheight}{18pt}
\setlength{\parskip}{5mm}

\lstset{language=Matlab,%
    basicstyle=\small,
    breaklines=true,
    morekeywords={matlab2tikz},
    keywordstyle=\color{blue},
    morekeywords=[2]{1}, keywordstyle=[2]{\color{black}},
    identifierstyle=\color{black},
    stringstyle=\color{mylilas},
    commentstyle=\color{mygreen},
    showstringspaces=false,
    numbers=left,
    numberstyle={\tiny \color{black}},
    numbersep=9pt, 
    emph=[1]{for,end,break},emphstyle=[1]\color{red},  
    xleftmargin=.2\textwidth, 
    xrightmargin=.2\textwidth,
    morecomment=[f][\color{Custom_color_1}][0]{\%},
}

\renewcommand{\rmdefault}{phv} 
\renewcommand{\sfdefault}{phv} 
\renewcommand{\algorithmcfname}{Algoritmo}
\pagestyle{fancy}

 
 

\begin{document}



\fancyhf{}
\renewcommand{\headrulewidth}{0pt}

\Huge
\begin{center}
\begin{tabular}{ll}
\raisebox{-.4\height}{\includegraphics[scale=0.15]{Figures/Logo_ipn.png}} & \textbf{ Instituto Politécnico Nacional}\\
\end{tabular}


\huge
\textbf{Unidad Profesional Interdisciplinaria en }  \\ 
\textbf{Ingeniería y Tecnologías Avanzadas} \\

\vspace{1cm}
\textbf{Unidad de Aprendizaje}

\rule{\textwidth}{1.5pt}

\textbf{Práctica $n$: Título de la práctica a realizar }
%\textbf{Reporte $n$: Título del reporte a realizar }

\rule{\textwidth}{1.5pt}

\vspace{1cm}
\Large

\textbf{Realizado por:}

\begin{center}
\begin{tabular}{| l | l | l |}
 \hline
 \textbf{Boleta} & \textbf{Nombre} & \textbf{Apellidos} \\
 \hline
 Integrante 1 & Integrante 1 & Integrante 1 \\ 
 \hline
 Integrante 2 & Integrante 2 & Integrante 2 \\ 
 \hline
 Integrante 3 & Integrante 3 & Integrante 3 \\ 
 \hline
\end{tabular}
\end{center}

\vspace{1cm}

\textbf{\today}

\end{center}


\newpage


\fancyhead{}
\fancyhead[C]{ \Large{Unidad de Aprendizaje}}
\fancyhead[L]{ \Large{UPIITA}}
\fancyhead[R]{ \Large{Grupo}}
\fancyfoot[R]{\thepage}
\renewcommand{\headrulewidth}{2pt}
\renewcommand{\footrulewidth}{2pt}
\pagenumbering{arabic}
\normalsize

\section{Introducción}

Descripción breve acerca del contenido y motivación de la práctica.

\section{Marco Teórico / Principios de operación}

En esta sección se deben incluir tanto la teoría que se utilizará descrita de una manera breve y concisa, así como las expresiones matemáticas requeridas. El texto de la parte teórica debe ser escrito por el estudiante de manera directa y no realizar copias textuales de internet. Las expresiones matemáticas que se utilizarán deben declararse desde el entorno de Latex y no utilizar impresiones de pantalla u otra clase de imágenes.

Las expresiones matemáticas son definidas como:
\begin{equation}
    x_{1}=2 x_{2}+8x_{3} +x_{4}^{2}
\end{equation}

En estas se pueden declarar tanto notación matemática como variables basadas en otros alfabetos (griego) tal que:
\begin{equation}
    z=\int_{0}^{t} \alpha(\tau) d\tau \qquad \alpha \in [0,\infty] \quad \alpha=0 \; \forall \; t \notin \mathbb{Z}
\end{equation}

Aquellas expresiones que contengan diversas ecuaciones en si mismas se escribe como:
\begin{equation}
    \begin{array}{rcl}
         z&=  & \beta q +2 r \\
         q&\neq  & \gamma z  \left( \sum_{i=1}^{n} r_{i} \right)
    \end{array}
\end{equation}

En caso de que el espacio entre las expresiones no sea el adecuado como se puede ver a continuación,
\begin{equation}
    \begin{array}{rcl}
         z&=  & \beta q +2 \frac{r}{3} \\
         q&\neq  & \gamma z  \left( \sum_{i=1}^{n} r_{i} \right)
    \end{array}
\end{equation}
\noindent se utiliza la siguiente estructura para ajustarlo,
{\setstretch{1.5}
\begin{equation}
    \begin{array}{rcl}
         z&=  & \beta q +2 \frac{r}{3} \\
         q&\neq  & \gamma z  \left( \sum_{i=1}^{n} r_{i} \right)
    \end{array}
\end{equation}}
Si se desea que tanto la fracción como la sumatoria tengan el tamaño apropiado la ecuación se reescribe como:
{\setstretch{1.5}
\begin{equation}
    \begin{array}{rcl}
         z&=  & \beta q +2 \displaystyle\frac{r}{3} \\
         q&\neq  & \gamma z  \left( \displaystyle\sum_{i=1}^{n} r_{i} \right)
    \end{array}
\end{equation}}

\subsection{Citar}

Para citar un documento tal que este aparezca como \cite{watson53} se utiliza el documento \textit{Bibliografia.bib}, donde se introducen los datos requeridos de cada una de las referencias como pueden ser: libros \cite{proakis1988introduction}, artículos científicos \cite{xie2019experimental}, trabajos en congresos \cite{wang2020orthogonal} y reportes técnicos \cite{world2019medication}.

Cuando se presentan diferentes fuentes de información respecto a un tema y se quieren englobar estas se puede implementar la citación de la siguiente manera \cite{guckenheimer2002chaos,foster1993significance,wang2007analysis,shorten2000hodgkin}. El orden con que se añaden las referencias al documento \textit{.bib} carece de importancia, Latex acomoda las referencias de forma automática respecto al orden en que son requeridas en el documento.

\section{Condiciones de experimentación}

En esta sección se debe de incluir lista de materiales utilizados durante la práctica en un formato de tabla como se presenta:

\begin{table}[ht]
\begin{center}
{
\setlength\arrayrulewidth{1pt}
\begin{tabular}{|l|c|}
\hline
\rowcolor{Custom_color_1} Material & Cantidad \\
\hline
\hline
Resistencia de $10k\Omega$ & $5$  \\
\hline
Resistencia de $1k\Omega$ & $10$   \\
\hline
Capacitor de $22pF$ & $2$  \\
\hline
PLA & $300g$   \\
\hline
Motor de DC de $8Nm$ a $12 V$ y $5A$ &  1\\
\hline
\end{tabular}
}
\end{center}
\caption{Lista de materiales.}
\label{Tab:Lista_materiales}
\end{table}

Se debe de incluir los diagramas requeridos ya sean electrónicos, mecánicos, etc. Estos deben ser realizados mediante el software correspondiente para que sean claramente entendibles y deben corresponde a diagramas y no a impresiones de pantalla de lo utilizado en entornos de simulación.

\begin{figure}[H]
    \centering
    \includegraphics{Figures/esquemaadquisicionori.eps}
    \caption{Ejemplo de diagrama.}
    \label{fig:Diagrama}
\end{figure}

También en dado caso que sea necesario se debe incluir el código utilizado, este debe presentarse en forma de pseudo-código para presentar el proceso realizado. No se debe incluir el código de manera directa a menos que se indique lo contrario. La manera en que se presentan estos resultados se puede observar en el Algoritmo \ref{al:PSO}.

\begin{algorithm}[H]
\SetAlgoLined
\KwData{$J(\cdot)$ and $\Omega_s$}
\KwResult{$x_{*}$}
Generar soluciones aleatorias $x_i(1)\in \Omega_s \quad i=1,\dots,r$ \;
Evaluar las soluciones $x_{i}(1)$ usando $J(\cdot)$ \;
Obtener los términos $x_{i,*}(1)$ and $x_{*}(1)$\;
\While{$k<k_{\text{max}}$}{
Evaluar las soluciones $x_i(k)$ usando $J(\cdot)$ \;
\eIf{$J\left( x_{*}(k) \right) \geq \alpha$}{
Calcular las velocidad $\omega_i(k)$ y las posiciones $x_i(k)$  \;
}{
\textbf{Detenerse}\;
}
$k=k+1$ \;
}
%$x^{*}=x_{*}$
\caption{Optimización por Enjambre de Partículas (PSO)}
\label{al:PSO}
\end{algorithm}

\section{Resultados}

Se deben incluir tablas con resultados, fotografías y gráficas según corresponda al trabajo a realizar. Todos y cada uno de estos elementos deben ser mencionados en el texto y referenciados adecuadamente. En los resultados se debe describir lo que se obtiene de forma clara y directa realizando comparaciones entre diversos resultados cuando esto sea conveniente.

``El estudio de los gatos que se realizó en este trabajo incluyó la realización de mediciones de peso y longitud de los felinos las cuales se incluyen en la Tabla \ref{Tab:Ejemplo}. Puede notarse que el peso medio de los sujetos de este estudio es de aproximadamente 4 kilogramos, mientras que las longitudes tienen su medio alrededor de los 40 centímetros. ''

\begin{table}[H]
\begin{center}
{
\setlength\arrayrulewidth{1pt}
\begin{tabular}{|l|c|c|}
\hline
\rowcolor{Custom_color_1} Nombre & Peso (Kg) & Logitud (centímetros) \\
\hline
\hline
Bigotes & $4.256$ & $35$  \\
\hline
Copo de nieve & $5.125$  & $45$   \\
\hline
Misifus & $3.5877$ & $37$  \\
\hline
Pazuzu & $4.058$  & $40$  \\
\hline
\end{tabular}
}
\end{center}
\caption{Tabla de ejemplo.}
\label{Tab:Ejemplo}
\end{table}

``El primer sujeto del estudio que se esta evaluando es \textbf{Bigotes}, este se presenta en la Figura \ref{fig:Gato} en la cual se puede notar que pertenece a la raza Siamés y esta acostumbrado a salir a exteriores. Este comportamiento ha provocado que presente múltiples infecciones debido a ... ''

\begin{figure}[H]
    \centering
    \includegraphics[scale=0.8]{Figures/gato.jpeg}
    \caption{Ejemplo de fotografía}
    \label{fig:Gato}
\end{figure}

\subsection{Formato de gráficas}

Las gráficas que se añadan deben de poseer título, definición de los ejes y leyenda donde se especifique a que corresponde cada una de las lineas de la gráfica con un tamaño de letra tal que sea posible visualizar fácilmente estos. Cada uno de los ejes debe poseer la numeración conveniente dado el caso y las dimensiones de los números deben ser adecuadas tal que sea claramente visible.
\begin{figure}[H]
    \centering
    \includegraphics[scale=0.8]{Figures/Grafica.eps}
    \caption{Ejemplo de gráfica}
    \label{fig:Grafica}
\end{figure}

Se incluye un ejemplo del código de Matlab donde se muestra como se realizó la gráfica presentada.

\lstinputlisting{Figures/myfun.m}

\section{Conclusiones}

Las conclusiones que se presenten deben se realizadas de manera individual por cada uno de los integrantes. Estas deben ser breves y concisas y referirse exclusivamente a los resultados obtenidos y las aseveraciones que pueden inferirse de ellos.

\subsection{Nombre1}

\lipsum[1]

\subsection{Nombre2}

\lipsum[2]

\subsection{Nombre3}

\lipsum[3]

\subsection{Nombre4}

\lipsum[4]

\subsection{Nombre5}

\lipsum[5]

\bibliographystyle{unsrt}
\small
\bibliography{Bibliografia}

\end{document}